\begin{result}{F.3.35}{The fourfold numerical quotient and its actual correspondence scope}{Known source theorem; scoped consequences proved}
Let \(k\) be a field, let \(\ell\) be a prime invertible in \(k\),
and use rational algebraic cycles and geometric \(\ell\)-adic cohomology.
For every abelian variety \(A/k\) of dimension at most four,
the identity on cycles gives an isomorphism
\[
 \operatorname{CH}^r(A)_{\Q}/{\sim_\ell}
   \xrightarrow{\ \sim\ }
 \operatorname{CH}^r(A)_{\Q}/{\sim_{\rm num}}
                         \qquad(r\in\Z).
 \tag{F.3.35a}
\]
In the contravariant convention for pure motives, the numerical
quotient is therefore fully faithful on the additive categories
with split idempotents generated by \(h(A)(m)\), where
\(\dim A\leq2\) and \(m\in\Z\). Its image contains every
numerical direct summand of a finite direct sum of those generators.
For a finite field \(k\) and a fourfold \(A/k\), the cycle map
\[
 \operatorname{CH}^2(A)\otimes\Q_\ell
    \longrightarrow H^4_{\acute{e}t}
           (A_{\overline k},\Q_\ell(2))^{G_k}
 \tag{F.3.35b}
\]
is surjective. These assertions retain the source's fourfold
dimension bound; the general numerical-quotient obligation in
F.3.5 and the fixed Betti comparison in E.2 remain on the page.
\end{result}
\begin{proof}
The surjectivity in (F.3.35b) is the exact statement of
\cite[Theorem 1.1, p.1]{RootBroe2026Fourfold}.
For a fourfold over an arbitrary field, (F.3.35a) is
\cite[Corollaries 1.2 and 8.2, pp.2 and 20]{RootBroe2026Fourfold}.
We opened the original August2026 preprint, read these statements,
the specialization map in Proposition6.2 on p.17, and the proof of
Corollary8.2 on p.20. The latter uses the specialization implication
of \cite[Corollary 3.19, p.13]{RootAncona2020Fourfold}:
the primitive algebraic cohomology of the generic fourfold embeds
in that of its good special fibre, compatibly with the intersection
pairing. Positivity on the special fibre forces the kernel of its
numerical quotient on this subspace to be zero. The cited fourfold
theorem and that exact implication are the external inputs here.

Here is the explicit transfer to lower dimensions. Choose an
elliptic curve \(E/k\), and, for \(d=\dim A<4\), form
\(B=A\times E^{4-d}\), its projection \(p:B\to A\), and the
zero section \(s:A\to B\). Such an elliptic curve exists over
every field: \(y^2=x^3-x\) works in characteristic different
from two, and \(y^2+y=x^3\) works in characteristic two.
For a numerically zero cycle \(z\) on \(A\) and any cycle
\(w\) of complementary dimension on \(B\), the projection formula
gives
\[
       \deg(p^*z\cdot w)=\deg(z\cdot p_*w)=0.
\]
Thus \(p^*z\) is numerically zero on the fourfold \(B\).
The source theorem gives \(\operatorname{cl}_\ell(p^*z)=0\).
Since \(p\circ s=\operatorname{id}_A\), functorial pullback gives
\(\operatorname{cl}_\ell(z)=s^*\operatorname{cl}_\ell(p^*z)=0\).
The reverse implication follows from the cycle-class product
and trace formulas as in F.3.5. This proves (F.3.35a) in all the
stated dimensions, including dimension zero and zero cycle groups.

For two stated generators the actual correspondence map is
\[
\begin{aligned}
 \operatorname{Hom}_{\ell}(h(A)(m),h(B)(n))
   &=\operatorname{CH}^{\dim A+n-m}(A\times B)_{\Q}/{\sim_\ell},\\
 \operatorname{Hom}_{\rm num}(h(A)(m),h(B)(n))
   &=\operatorname{CH}^{\dim A+n-m}(A\times B)_{\Q}/{\sim_{\rm num}}.
\end{aligned}
\tag{F.3.35c}
\]
Their product has dimension at most four, so (F.3.35a) identifies
these two spaces by the identity on cycles. It preserves composition
because the same pullback, product, and proper pushforward formula
defines composition before and after the quotient. On finite direct
sums the maps are the corresponding matrix maps. For projectors
\(e,f\), restrict the isomorphism to \(f\operatorname{Hom}e\).
Finally a numerical idempotent lifts uniquely through the endomorphism
isomorphism, and its lift is idempotent by injectivity. These are
the asserted fully faithful functor and its stated image.

The scope needed for larger motivic categories can now be read from
(F.3.35c). For instance, endomorphisms of \(h(A)\) for a fourfold
use cycles on \(A\times A\), of dimension eight; arbitrary tensor
powers use the corresponding still larger products. No equality
on these products is inserted into (F.3.35a). Likewise, the cycle
specialization in the source sends a cycle on a chosen characteristic
zero lift to its special fibre. Its explicit map is specialization
of the closure of that cycle, followed by the compatible smooth
proper cohomology map. Applying it to a finite-field reduction
does not construct a cycle on the originally specified number-field
fourfold. That lifting map and the tensor-invariant comparison in
its fixed Betti realization remain to be supplied for the proposed
application to E.2.
\end{proof}
